\section*{Abstract}
Doc.~287 classifies FFGFT's parallel representations by relationship type (identity, bijection,
projection). This document reads the same representations \emph{differently}: not as a map but as a
\emph{toolbox}. The signal-engineering experience -- transforms exist because they make computations
cheap (convolution becomes multiplication) -- carries over directly. The mass operator is a
Z$_3$-circulant; the DFT diagonalises it. From this follow five concrete, machine-precision-checked
computational levers: (1)~the lepton mass roots are a \emph{3-point DFT} instead of an eigenvalue
problem; (2)~Koide $Q=2/3$ falls out in \emph{two Parseval lines} and is manifestly phase-independent;
(3)~\emph{any} operator function (power, inverse, exponential) is a scalar evaluation at three
eigenvalues; (4)~chaining Z$_3$ couplings is \emph{pointwise multiplication} in the DFT domain;
(5)~the $\xi$ recursion is a \emph{first-order IIR filter} with a pole at $1-\xi$, and the three
regimes are pole locations. Honestly: this is computational leverage, not a derivation -- the
transform moves $\sqrt2$ and $\theta$, it does not create them.

\section{The idea: changing representation as a computational lever}
In signal engineering one changes representation not by taste but for computational gain. A convolution
in the time domain is a pointwise multiplication in the frequency domain; a linear time-invariant
system is an algebraic transfer function in the $z$ domain; magnitude and phase separate, so they are
computed independently. Exactly this leverage sits in FFGFT's parallel representations -- one only has
to read them as transforms rather than worldviews.

The key is an observation from Doc.~286: the mass operator \emph{is} a Z$_3$-circulant. A circulant is
diagonalised by the discrete Fourier transform (DFT) -- the signal-engineering principle \enquote{cyclic
convolution $\leftrightarrow$ DFT}. Everything else follows.

\begin{keyresult}
The Z$_3$-circulant mass operator $C$ is diagonalised by the unitary DFT matrix $F$:
$C = F^{\dagger}\Lambda F$, $\Lambda=\mathrm{diag}(\lambda_0,\lambda_1,\lambda_2)$. Thus every
\enquote{function of $C$} problem becomes an evaluation at the three scalars $\lambda_k$.
\end{keyresult}

\section{The setup: first row and DFT}
The Foot-Koide configuration $\sqrt{m_k}=1+\sqrt2\cos(\theta+2\pi k/3)$ is the DFT of a \emph{3-vector},
the first row of the circulant:
\[
  c=\Bigl(1,\ \tfrac{\sqrt2}{2}\,e^{i\theta},\ \tfrac{\sqrt2}{2}\,e^{-i\theta}\Bigr),
  \qquad r=\sqrt2,\quad \theta=\tfrac{2}{9}.
\]
Here the magnitude $r=\sqrt2$ sits in $|c_1|=|c_2|=\sqrt2/2$ and the phase $\theta$ in the argument --
magnitude and phase parts are separated from the start. $C$ is Hermitian (checked:
$\max|C-C^{\dagger}|=0$), hence real-spectral.

\section{Lever 1: eigenvalues without a characteristic polynomial}
The three mass roots $\sqrt{m_k}$ are the DFT of $c$: $\lambda_k=\sum_{j}c_j\,\omega^{jk}$,
$\omega=e^{2\pi i/3}$. No $3\times3$ eigensolver, no polynomial -- a DFT. Numerically
$\sqrt m=(0{.}0404,\ 0{.}5802,\ 2{.}3794)$, matching the closed Foot-Koide form and \texttt{eigvalsh}
to $8{.}9\cdot10^{-16}$ (\texttt{dft\_eigenwerte.py}).

\section{Lever 2: Koide Q in two Parseval lines}
With $\lambda_k$ as DFT coefficients, Parseval gives directly:
\[
  \sum_k \sqrt{m_k}=3\,c_0=3,\qquad \sum_k m_k=3\sum_j|c_j|^2=6,
  \qquad Q=\frac{\sum m_k}{(\sum\sqrt{m_k})^2}=\frac{6}{9}=\frac{2}{3}.
\]
$\sum\sqrt m$ reads only the DC term $c_0$, $\sum m$ only the energy $\sum|c_j|^2$. Both depend
\emph{only on the magnitude} $|c_j|$, i.e.\ on $r$ -- not on $\theta$. Over all $\theta\in[0,2\pi)$,
$Q$ stays exactly $0{.}666667$; a different magnitude shifts $Q$ ($r=1\!\to\!1/2$,
$r=\sqrt2\!\to\!2/3$, $r=1{.}6\!\to\!0{.}76$). Thus the radial/angular conclusion of Doc.~282 is not a
separate statement but a trivial Parseval consequence: the magnitude spectrum fixes the invariant, the
phase spectrum does not enter (\texttt{parseval\_koide.py}).

\section{Lever 3: any operator function for free}
From $C=F^{\dagger}\Lambda F$ follows $f(C)=F^{\dagger}f(\Lambda)F$ for any function $f$. An operator
function is therefore three scalar evaluations $f(\lambda_k)$ instead of matrix algebra:
\begin{itemize}
  \item \textbf{power} $C^n$ -- for recursion / scale flow;
  \item \textbf{inverse} $C^{-1}$;
  \item \textbf{exponential} $\exp(C)$ -- for time evolution / propagator.
\end{itemize}
Checked against \texttt{matrix\_power}, \texttt{inv} and a Taylor series: $C^5$ to
$3{.}6\cdot10^{-14}$, $C^{-1}$ to $4{.}8\cdot10^{-15}$, $\exp(C)$ to $2{.}7\cdot10^{-15}$
(\texttt{operatorfunktionen.py}). This is exactly the sector the dynamic part (Doc.~286) heads into:
there one needs functions of the operator, and they are trivial in the DFT domain.

\section{Lever 4: convolution becomes multiplication}
A Z$_3$-cyclic coupling between the three generations is a circular convolution. The convolution
theorem turns it, in the DFT domain, into a pointwise multiplication:
$\mathrm{DFT}(a\ast b)=\mathrm{DFT}(a)\cdot\mathrm{DFT}(b)$. Chaining two Z$_3$ couplings -- a matrix
product of two circulants -- becomes three scalar products. Checked to $9\cdot10^{-16}$
(\texttt{faltung\_multiplikation.py}). This is the actual signal-engineering core: the expensive
operation (matrix product) becomes cheap (three multiplications) once one computes in the right
representation.

\section{Lever 5: the recursion is a filter}
$r(k{+}1)=r(k)(1-\xi)$ is a first-order IIR filter with $z$-transfer function
$H(z)=1/(1-(1-\xi)z^{-1})$, a pole at $z=1-\xi=0{.}99986667$. $|z|<1$ means stable / contracting; the
pole's distance to the unit circle is exactly $\xi$, hence the time constant $k^{*}\sim1/\xi=7500$
steps. The three regimes of Doc.~285/286 are pole locations: inside (negative feedback, contracting),
on the circle (Barkhausen limit), outside (positive feedback, run-away). Impulse response and direct
recursion agree to $10^{-15}$ (\texttt{rekursion\_iir.py}). Log-periodicity, the $\xi\!\to\!\varphi$
passage and decompactification-as-runaway are thus pole computations, not a separate formalism.

\section{Magnitude and phase, minimum phase and all-pass}
The five levers share a common ground that repeats the θ work (Doc.~286): in the spectral
representation, magnitude and phase separate. Everything FFGFT determines is magnitude content ($\xi$,
hierarchy, $Q=2/3$, $\langle f\rangle$); the one open value $\theta$ is pure phase. In signal terms
FFGFT's real recursion is a minimum-phase / amplifier stage (the magnitude carries everything), and a
phase stage $e^{i\varphi}$ is an all-pass (magnitude-preserving) -- exactly HLV's skew-adjoint $R$
(BD17A). The all-pass can be \emph{appended} without touching the magnitude (and hence the masses and
$Q$); this is why the phase is computationally separable and the open remainder.

\section{What this means in practice}
The gain is wherever one needs \emph{functions of the operator} or \emph{chains of couplings} -- i.e.\
in the dynamic / kinetic sector (Doc.~286), at scale flow and time evolution. Instead of $3\times3$
matrix algebra one computes at three eigenvalues; instead of matrix multiplication three scalar
products; instead of recursion loops one pole location. The radial invariants (masses, $Q$) are
Parseval one-liners. The toolbox is standard signal engineering applied to the Z$_3$ core.

\section{Honest boundary}
\begin{important}
This is computational leverage, \emph{not} new physics (P35). The change of representation \emph{moves}
$\sqrt2$ and $\theta$, it does not \emph{create} them -- both remain inputs. The full leverage is on
the C$_3$/Z$_3$ sector (3$\times$3 circulant); the full $T^4$ with the $\varphi$ harmonics is larger,
where the DFT diagonalisation is only the generation block. And no lever changes the status of the open
phase: it only makes visible \emph{why} it stands separately.
\end{important}

\begin{conclusionBox}[Balance: the representations compute too]
The parallel representations are not only classification (Doc.~287) but computational tools. Because
the mass operator is a Z$_3$-circulant, the DFT diagonalises it, and five standard signal-engineering
levers become applicable: eigenvalues as DFT, invariants as Parseval, operator functions as scalar
evaluation, chaining as multiplication, recursion as filter pole. All five checked to machine precision.
The gain is concrete in the dynamic sector; the boundary is honest: a faster computational route, not an
additional derivation. The signal-engineering view missing in Doc.~287 is thereby added as its own
toolbox.
\end{conclusionBox}
